\section{Discussion and Conclusion}
The prevailing assumption in cybersecurity architecture is that Deep Learning and rigorous interpretability are mutually exclusive within the rigid confines of Edge Computing. This research empirically invalidates that dichotomy. By deploying the \textbf{X-IDS} framework, we have demonstrated that mathematically rigorous Explainable AI (XAI) is not merely a theoretical post-hoc luxury, but a viable, highly resilient operational tool for resource-constrained environments.

\subsection{Operational Viability and SOC Benefits}
Through the integration of Apache Kafka as a localized buffer and PySpark for micro-batch execution, X-IDS successfully maintained state-of-the-art anomaly detection (F1-score of $99.18\%$) directly on a standard 4GB Raspberry Pi 4. Crucially, the system proved that utilizing SHAP (SHapley Additive exPlanations) works effectively as a deterministic dimensionality reduction strategy. Shrinking the data throughput mathematically insulated the Raspberry Pi from memory exhaustion. 

From an operational perspective, the value of X-IDS resides in its capacity to generate trust. Traditional binary IDS alerts trigger a cascaded investigative burden upon Security Operations Center (SOC) analysts, demanding hours of manual packet analysis. Conversely, by accepting a microscopic sub-30ms computational overhead, X-IDS provides analysts with real-time SHAP matrices that definitively explain \textit{why} an alert was generated. Pinpointing precise anomalies, such as an anomalous variance in \texttt{Bwd Packet Length}, empowers incident responders to isolate threats instantaneously with mathematical certainty.

\subsection{Architectural Limitations}
While X-IDS establishes a robust baseline for transparent Edge security, empirical benchmarking revealed absolute hardware limitations. During simulated volumetric Distributed Denial of Service (DDoS) attacks exceeding $15,000$ flows per second, the Raspberry Pi's CPU utilization sustained an unstable $98\%$. Even with the SHAP-pruned dataset, the localized Spark JVM and the Kafka consumer struggled to share the processor's limited L2 cache, resulting in localized Kafka partition lag. This indicates that while the Raspberry Pi excels at analyzing standard IoT telemetery pipelines, it cannot operate as a standalone sink for sustained, infrastructure-level DDoS deluges without additional computational relief.

\subsection{Future Work}
To address these physical hardware constraints and further optimize the intersection of AI and Edge computing, future research will explore three primary vectors:
\begin{enumerate}
    \item \textbf{Hardware Acceleration:} We intend to migrate the inference engine from standard ARM CPUs to specialized Edge AI accelerators, such as the Google Coral Edge TPU or NVIDIA Jetson Nano. This will fundamentally resolve the computational bottlenecks of real-time SHAP generation via parallelized tensor processing.
    \item \textbf{Model Quantization:} Future iterations will explore aggressive post-training quantization techniques (e.g., INT8 precision) to further compress the LightGBM models, theoretically doubling inference throughput without sacrificing decision accuracy.
    \item \textbf{Federated Learning (FL):} To preserve zero-trust privacy architectures, we will extend X-IDS into a Federated Learning topology. This will allow highly distributed, resource-constrained Edge nodes to collaboratively train a global IDS model localized telemetry, entirely eliminating the need to expose raw network data to centralized cloud clusters.
\end{enumerate}
